\section{Conclusion}
\label{sec:conclusion}

This study provides a comprehensive characterization of the localized magnetic environments at candidate lunar landing sites and evaluates their potential impact on biological systems. Our analysis demonstrates that while the lunar crust exhibits extreme magnetic heterogeneity, all prospective Artemis, Apollo, and Chang'e landing sites constitute severe hypomagnetic field environments, with fields typically less than \SI{5}{\nano\tesla} at \SI{30}{\kilo\meter} altitude. Synthesis of current magnetobiology literature suggests that such sustained exposure to hypomagnetic conditions could profoundly alter plant development, microbial ecology, and potentially human cellular repair mechanisms. As the Artemis program advances toward sustained lunar surface operations, the profound absence of a normative biological magnetic field must be integrated into risk assessments. We strongly recommend the inclusion of high-resolution surface magnetometry and controlled magnetic-biology experiments in early lunar missions to empirically delineate the isolated physiological effects of the lunar hypomagnetic environment.
